\section{Discussion}\label{sec:discussion}

The purpose of this benchmark is not to replace the practical pipeline of the blinded precursor, but to determine which conclusions remain when the temporal split, optimizer budget, validation objective, and economic endpoint are made explicit. The resulting evidence separates a model's average test-block outcome, its sensitivity to the optimizer--seed block, and its simplified trading consequence. This distinction matters for intraday financial prediction, where a change in selection criterion can alter a model ranking without demonstrating a stable change in market-wide forecasting ability.

The present design preserves the precursor's technical-indicator workflow and fixed seven-feature representation, but changes the comparison unit. Rather than reporting a GA-selected configuration after randomized cross-validation, every backbone is evaluated with RS, GA, PSO, DE, and GWO under the same 1,000-candidate allowance and the same chronological train/validation/test partition. Random Search is therefore a budget-matched reference against which directed population methods can be interpreted \cite{bergstra2012random_search}. The test features themselves are not used during candidate selection. However, because the target is shifted by one retained row, one training target crosses into validation and two validation targets cross into the nominal test interval. The test is therefore nominally subsequent but not strictly isolated at these three labels, and conclusions are interpreted as exploratory pending a purged rerun.

This control improves internal comparability, but it narrows the claim that the data can support. The evidence concerns four classifier families on one five-minute WIN period, conditional on an inherited input representation. It does not show that the seven indicators would be selected again in another training period, nor that an optimizer with a favourable result here will retain that advantage under another market regime. The contribution is consequently a controlled experimental frame, not a universal ordering of metaheuristics.


The test-block summaries in Table~\ref{tab:predictive_summary} show a clear reversal at the model level. Under MCC/$F_1$ fitness, RF has the largest marginal means for accuracy, MCC, and $F_1$, and its small standard deviations indicate comparatively low sensitivity to the evaluated optimizer--seed blocks. Under weighted-accuracy fitness, MLP is first on all three reported predictive measures, while 1D-CNN is also stable but remains below MLP in mean accuracy and MCC. The change in ordering is descriptive. Because the two protocols differ in both the selection metric and the inclusion of in-sample training accuracy, it cannot be attributed to metric choice alone.

MCC is useful because it summarizes the association between predicted and observed binary classes beyond raw accuracy \cite{chicco2020mcc}; however, optimizing an imbalance-aware criterion cannot guarantee the best outcome for every backbone and search trajectory. Conversely, the MLP advantage under weighted-accuracy fitness is concentrated in a low-variance regime, whereas MLP, SVM, and 1D-CNN outcomes under MCC/$F_1$ fitness vary much more across the evaluated blocks. A reported mean must therefore be accompanied by its dispersion and by the selection rule that generated it.

The rank analysis reinforces this caution. Under MCC/$F_1$ fitness, RF has the highest marginal mean MCC, but MLP has the lowest mean Friedman rank (1.533) across the matched blocks. This is not a contradiction: means weight score magnitude, whereas ranks summarize ordering within each block. Under weighted-accuracy fitness, both summaries place MLP first, and its Holm-adjusted contrasts against all other backbones are significant. Because all 15 blocks share one chronological test segment, these results indicate separation within the evaluated design rather than independent replication across time.


Figures~\ref{fig:rf_convergence_mccf1}--\ref{fig:cnn_convergence_accuracy} index progress by candidate evaluations rather than by optimizer-specific generation counts. Each line shows the mean trajectory over the three common seeds; variability envelopes are omitted to keep closely spaced optimizer curves legible. In each panel, the main y-axis is tightly bounded by the observed range of the five mean trajectories, with a small margin, and an inset magnifies the final 250 evaluations. Because limits are panel-specific, visual heights should not be compared across panels without reading their y-axis scales. Early gains indicate how quickly an optimizer locates useful regions of a search space, whereas later plateaus show the incremental benefit of additional evaluations under the fixed budget. Convergence in validation fitness does not establish equivalent generalization, so the test-block values and their dispersion in Table~\ref{tab:predictive_summary} remain necessary for interpretation.

Search behaviour interacts with the backbone and the objective. A method that improves early for RF need not be the most reliable procedure for MLP or 1D-CNN, whose mixed discrete--continuous hyperparameter spaces have different response surfaces. The article therefore does not claim a universal winner among RS, GA, PSO, DE, and GWO. A stronger optimizer-level claim would require a dedicated analysis across more seeds and temporal replications, with the optimizer itself as the inferential factor.


Table~\ref{tab:economic_summary} makes the distinction between predictive evaluation and the exploratory trading rule concrete. With next-bar execution, 1D-CNN has the largest mean net profit under both fitness protocols. Under MCC/$F_1$ fitness, RF has the least negative mean drawdown; under weighted-accuracy fitness, 1D-CNN leads mean profit, profit factor, and annualized daily Sharpe. Dispersion remains substantial for several model--protocol combinations, so the apparent return leader does not establish a stable economic advantage.

These values are reported as points under a fixed next-bar long/short rule with a five-point round-trip cost. They are not an investable-strategy claim: the evaluation does not model a full transaction-cost schedule, market impact, position sizing, or a live signal-selection policy. Its role is diagnostic: it shows that converting a one-bar-ahead regime prediction to trades can change the practical ranking of model families, and that economic reporting should accompany predictive metrics.

\subsection{Limitations and Research Agenda}

Seven limitations delimit the present evidence. First, only three seeds are common to every model--optimizer cell; standard deviations and rank tests describe the evaluated blocks, not repeated market samples. Second, exact library versions for the official TensorFlow and scikit-learn runs are not archived, and the precursor Information-Gain ranking can be traced only to the recorded seven-feature list and the researcher's recollection of a WEKA ranking on the January--March training segment; the original WEKA project, output, and settings are unavailable for independent replay. Third, the test covers one instrument and one chronological segment, so regime sensitivity and cross-market transfer remain unknown. Fourth, the Information-Gain mask was inherited from the blinded precursor rather than re-estimated on the current training segment; all conclusions are conditional on that representation. Fifth, labels are derived from retrospective hourly pivots and shifted to the next retained row. One training target crosses into validation and two validation targets cross into the nominal test block; these three rows were not purged. Four per-ticker terminal rows also lack a next-row target and retain their contemporaneous source label. A cutoff-only rerun found no other changed training or validation labels (9,032 and 3,008 checked, respectively), but this does not remove direct boundary crossings or establish causal availability of retrospective pivot labels. Consequently, test isolation is incomplete at these boundaries, and the reported comparisons should be treated as exploratory pending a purged rerun. Sixth, equalizing the budget at 1,000 evaluations establishes fairness at one computational scale but does not characterize larger-budget behaviour or wall-clock efficiency; the fit-count figure is a hardware-independent proxy and uses a run-relative threshold. Seventh, no-tuning baselines are available only for the majority rule and default RF, so the benefit of search is not isolated for SVM, MLP, or 1D-CNN.

The next study should preserve chronological evaluation and add a preregistered common 20-seed design across multiple non-overlapping market periods. It should retain pivot-confirmation timestamps, purge boundary samples whose target regime uses observations beyond the training or validation cutoff, and rerun selection and evaluation under that purged design. It should extend the present naive and default-RF references to no-tuning baselines for every backbone; archive every candidate evaluation; report compute time; and test optimizer differences separately from model differences. Re-estimating feature selection within each training window would further distinguish the effect of representation from hyperparameter search.
